% Source: https://artofproblemsolving.com/community/c3h2505412_what_does_this_usually_mean
% Stable identity: aops-topic-2505412-q01
In number theory, $\tau(n)$ usually denotes the number of positive divisors of the positive integer $n$.
\begin{enumerate}
\item[(a)] Prove that if $n=p_1^{e_1}p_2^{e_2}\cdots p_k^{e_k}$ with distinct primes $p_i$, then $\tau(n)=(e_1+1)(e_2+1)\cdots(e_k+1)$.
\item[(b)] Compute $\tau(2021)$ and $\tau(10!)$.
\item[(c)] Find the smallest positive integer $n$ with $\tau(n)=12$.
\end{enumerate}
